\chapter{Introduction}
\label{ch:intro}

\chapquote{``The best interface is no interface until you need one.''}{Anonymous UX aphorism}

\section{Motivation}
\label{sec:motivation}

Modern machine learning (ML) has produced spectacular headline results
in image recognition, natural-language processing, generative modelling,
and scientific discovery. However, the day-to-day experience of a
student, applied researcher, or first-year data engineer confronted with
a new tabular dataset is far less glamorous. Even relatively small
projects require the practitioner to (i)~load and profile the data,
(ii)~identify and encode categorical variables, (iii)~engineer missing
values, (iv)~pick a family of models, (v)~configure and search
hyperparameters, (vi)~compare candidates on a validation split with a
sensible metric, (vii)~inspect feature importance and residuals, and
(viii)~produce a reproducible record of what was done.
Each of these steps has well-established best practices, and each is
easy to get subtly wrong. Automated Machine Learning (AutoML) systems
\citep{feurer2015autosklearn, olson2016tpot, ledell2020h2o, jin2019autokeras}
were introduced to reduce the cognitive load of these repeated
decisions and to make sensible defaults available to non-experts.

Meanwhile, the recent generation of instruction-tuned large language
models (LLMs) --- exemplified by the GPT family
\citep{brown2020gpt3, openai2023gpt4}, the Llama family
\citep{touvron2023llama, llama3card}, and the Claude family --- has
demonstrated a striking ability to plan and to invoke external ``tools''
in response to natural-language requests
\citep{schick2023toolformer, yao2023react, patil2023gorilla,
openai2023function}. The Model Context Protocol (MCP)
\citep{mcpSpec2024, anthropic2024mcp} standardises how such tools are
discovered, described, and invoked, which makes it feasible to expose
an entire application's capabilities as an inspectable tool catalog
that an LLM can drive.

The intersection of these two trends is the design space explored by
this project: an AutoML platform whose \emph{interface} is a
conversation, whose \emph{implementation} is deterministic Python, and
whose \emph{contract with the LLM} is a well-typed MCP tool schema.

\section{Evolution of AutoML}
\label{sec:evolution-automl}

Automated machine learning has traversed several distinct generations:

\begin{itemize}[leftmargin=1.3em]
  \item \textbf{First generation --- hyperparameter search.} Grid search
        and random search over a small number of scikit-learn
        \citep{sklearn2011} hyperparameters, often driven by
        \code{GridSearchCV} in a notebook. Random search is a
        surprisingly strong baseline \citep{bergstra2012random}.
  \item \textbf{Second generation --- Bayesian and evolutionary
        pipeline search.} Systems such as Auto-WEKA
        \citep{thornton2013autoweka}, Auto-sklearn
        \citep{feurer2015autosklearn}, and TPOT \citep{olson2016tpot}
        automated pipeline construction using Bayesian optimisation
        (SMAC \citep{hutter2011smac}) or genetic programming, and
        popularised the idea of a ``leaderboard'' of pipelines.
  \item \textbf{Third generation --- boosting-first AutoML and neural
        architecture search.} H2O AutoML \citep{ledell2020h2o},
        AutoGluon \citep{erickson2020autogluon}, and Auto-Keras
        \citep{jin2019autokeras} placed gradient-boosted trees and
        stacked ensembles at the centre of tabular AutoML while NAS
        systems \citep{zoph2017nas, real2019regularized} attacked deep
        image and language models. Modern gradient boosters
        \citep{chen2016xgboost, ke2017lightgbm, prokhorenkova2018catboost}
        remain the state of the art on tabular data.
  \item \textbf{Fourth generation --- efficient hyperparameter
        optimisation.} Optuna \citep{akiba2019optuna} popularised
        define-by-run search spaces and TPE
        \citep{bergstra2011algorithms} sampling with pruning; Hyperband
        \citep{li2017hyperband} and BOHB \citep{falkner2018bohb}
        introduced multi-fidelity approaches.
  \item \textbf{Fifth generation --- agentic and conversational
        AutoML.} A very recent line of work
        \citep{hollmann2024caafe, guo2024dsagents} uses LLMs to plan
        and execute portions of the ML pipeline. The
        \emph{platform} described here belongs to this fifth
        generation.
\end{itemize}

\section{Problem Being Addressed}
\label{sec:problem-addressed}

Even with all of the above tools, an important gap remains:

\begin{quote}
\emph{There is no lightweight, single-user, open, reproducible platform
that lets a researcher upload a CSV, describe the intent in natural
language, and receive back a Jupyter notebook, a trained model, a set
of plots, and a PDF report --- with a clear, inspectable record of
every tool the system called, every argument it used, and every artefact
it produced.}
\end{quote}

Commercial vendors offer end-to-end products but hide their engineering
and require accounts, API keys, and network egress. Research systems
such as Auto-sklearn \citep{feurer2015autosklearn} and AutoGluon
\citep{erickson2020autogluon} are excellent SDKs but not products
--- they lack a chat interface and a reproducible artefact story.
LLM ``agents'' such as ChatGPT's code interpreter can execute Python
but their tool catalog is monolithic and not easily extensible by a
third party. \projname{} fills this gap.

\section{Why Conversational, Agentic AutoML is Useful}
\label{sec:why-agentic}

Three properties of conversational, tool-using AutoML are attractive
for the target user:

\begin{enumerate}[leftmargin=1.4em]
  \item \textbf{Low activation energy.} A student can upload a CSV and
        write ``run a full EDA'' without recalling scikit-learn API
        surface. The LLM handles the mapping from intent to tool
        invocation.
  \item \textbf{Auditable execution.} Unlike a black-box AutoML SaaS,
        every tool call, argument, and result is displayed inline and
        stored, so a reviewer can reproduce the run manually.
  \item \textbf{Cheap extensibility.} Adding a new capability (e.g.\ a
        new plot type, a new estimator, a new report format) is a
        matter of adding a decorated function to one of the MCP
        servers --- no orchestrator changes and no front-end changes
        are required.
\end{enumerate}

\section{Project Motivation}
\label{sec:project-motivation}

The specific motivation for building \projname{} was the observation
that \emph{first-year data-science students frequently produce
technically correct but conceptually shallow analyses} because they
do not know which questions to ask. A conversational tool that models
those questions --- ``what does the target column look like?'', ``are
any features strongly correlated?'', ``which model wins on this
dataset?'', ``why does it win?'' --- can, if designed carefully,
teach as it works.

A second motivation was that the Model Context Protocol
\citep{mcpSpec2024} had matured to a point where a rich, in-process
tool catalog could be built with a single Python dependency (FastMCP
\citep{fastmcp}), without any HTTP plumbing or authentication layers.
This made a research-grade prototype feasible for a single-author
project.

\section{Objectives}
\label{sec:objectives}

The concrete objectives of the project are:

\begin{enumerate}[leftmargin=1.4em]
  \item Design and implement a modular, in-process MCP server layout
        that covers the core stages of a tabular AutoML workflow ---
        ingestion, validation, EDA, training, explainability, export.
  \item Implement an LLM-driven orchestrator that unifies the tools
        into a single catalog and routes tool calls to the correct
        server, with support for streaming responses and per-tool
        progress updates.
  \item Provide a filesystem-based artefact registry that records every
        generated file with a stable identifier, a category, and
        arbitrary metadata.
  \item Build a thin Streamlit front-end that supports upload, chat,
        and artefact browsing while remaining a client of the same
        orchestration API that is used programmatically.
  \item Produce a project report that documents the architecture, the
        MCP integration, and the AutoML pipeline in enough detail to
        allow future contributors to extend the platform.
\end{enumerate}

\section{Contributions}
\label{sec:contributions}

The specific contributions of this work are:

\begin{itemize}[leftmargin=1.4em]
  \item An \emph{end-to-end, single-machine} AutoML platform built on
        MCP whose five services and eleven tools cover the tabular
        workflow from raw CSV to PDF report.
  \item An \emph{in-process transport} for the MCP client--server
        protocol that eliminates HTTP overhead while preserving the
        protocol's tool discovery, invocation, and progress semantics.
  \item A \emph{deterministic-first} design in which the LLM plans
        and formats but never computes: numerical results always
        originate from Python code.
  \item A \emph{filesystem artefact registry} that is trivial to
        version-control, inspect, and reset, and that decouples
        artefact generation from any storage vendor.
  \item A reference implementation, in $\approx$2\,600 lines of
        Python plus a Streamlit UI, that can serve as an educational
        substrate for further research on agentic data science.
\end{itemize}

\section{Scope}
\label{sec:scope}

This report is scoped to the current state of the repository
\file{automl-mcp-platform}, in which the platform:

\begin{itemize}[leftmargin=1.4em]
  \item Runs on a \textbf{single machine} for a \textbf{single user};
        multi-tenancy has been intentionally removed.
  \item Handles \textbf{tabular} data in CSV, TSV, or Parquet format;
        image, text, and multimodal data are out of scope.
  \item Uses the \textbf{Groq} Chat Completions API as the LLM backend;
        other providers are supported in principle via the standard
        OpenAI-compatible schema but are not tested.
  \item Persists artefacts on the \textbf{local filesystem} under
        \file{data/}; no cloud storage or database is used.
  \item Provides two entry points: a \textbf{Streamlit UI} for
        interactive use and a small \textbf{Python API}
        (\code{orchestrator.stream\_chat}) for programmatic use.
\end{itemize}

\section{Report Organization}
\label{sec:report-org}

The remainder of this report is organised as follows.
Chapter~\ref{ch:background} reviews the technical foundations of
AutoML, MCP, and LLM tool use.
Chapter~\ref{ch:litreview} surveys the academic literature relevant to
the design.
Chapter~\ref{ch:problem} formalises the problem and requirements.
Chapter~\ref{ch:requirements} enumerates functional and non-functional
requirements.
Chapter~\ref{ch:arch} presents the system architecture.
Chapter~\ref{ch:design} discusses the platform design in module-level
detail.
Chapter~\ref{ch:mcp} focuses specifically on the MCP integration and
the agent decision flow.
Chapter~\ref{ch:pipeline} traces the end-to-end AutoML pipeline.
Chapter~\ref{ch:impl} presents key implementation excerpts.
Chapter~\ref{ch:experiments} defines the experimental protocol and
Chapter~\ref{ch:results} discusses the results.
Chapter~\ref{ch:demo} maps a complete run onto a YouTube-style
demonstration.
Chapter~\ref{ch:limitations} states the honest limitations, and
Chapter~\ref{ch:conclusion} concludes and discusses future directions.
Appendices provide the project structure, configuration reference,
tool signatures, reproducibility notes, and additional results.
